\section{Simulated models: equations and parameters}\label{app:models}
For a synapse from population $Y$ to $X$, the resulting first-generation neural
mass model (NMM1) equations are
\cite{ruffini2025rosetta}
\begin{equation}
\hat L_{X\leftarrow Y}\big[u_{X\leftarrow Y}\big]=C_{X\leftarrow Y}\,r_Y,\qquad
v_X=\Lambda[E]+\!\!\sum_{Y:\,C_{X\leftarrow Y}\neq0}\!\!u_{X\leftarrow Y},\qquad
r_X=\Sig(v_X).
\label{eq:nmm1}
\end{equation}
The operator $\hat L$ is the inverse synaptic filter, $u_{X\leftarrow Y}$ its
postsynaptic potential, $C_{X\leftarrow Y}$ the connectivity, and $\Lambda[E]$
the direct polarization caused by the applied field $E$ (V/m). The rates $r_X,r_Y$
are in s$^{-1}$, the potentials $v_X,u_{X\leftarrow Y}$ in mV, and the
pathway weights $C_{X\leftarrow Y}$ dimensionless. This field contribution enters
alongside the synaptic potentials, without first passing through a synapse.

Synaptic operators convert rates into postsynaptic potentials, which sum with
the direct field polarization before the rate curve is applied.

\subsection*{Jansen--Rit column}
Writing $y_0$ for the excitatory PSP driving the interneurons, $y_1$ and $y_2$ for the
excitatory and inhibitory PSPs onto the pyramidal population, external input $p$, and the
applied-field perturbation $s(t)$ (\S\ref{sec:coupling}) entering the pyramidal sigmoid,
\begin{align}
\ddot y_0 &= Aa\,\Sig\!\big(y_1-y_2+s(t)\big)-2a\dot y_0-a^2 y_0, \label{eq:jr0}\\
\ddot y_1 &= Aa\,\big[\,p+C_2\,\Sig(C_1 y_0)\,\big]-2a\dot y_1-a^2 y_1,\\
\ddot y_2 &= Bb\,C_4\,\Sig(C_3 y_0)-2b\dot y_2-b^2 y_2,
\end{align}
The pyramidal output passes through $h_e$ to give $y_0$; the excitatory
and inhibitory interneurons receive $C_1y_0$ and $C_3y_0$, and return through
weights $C_2$ and $C_4$. The excitatory and inhibitory impulse responses are
$h_e(t)=Aa\,t e^{-at}$ and $h_i(t)=Bb\,t e^{-bt}$ for $t\ge0$.
The equations are integrated as six first-order ODEs; the LFP proxy is $v=y_1-y_2$ \cite{jansen1995}.
Table~\ref{tab:params} gives the reference column parameters. The synaptic rates
correspond to excitatory and inhibitory decay times of $10$ and $20$~ms, while
$A$ and $B$ set the postsynaptic-potential gains. The sigmoid parameters fix its
voltage scale and rate ceiling; the sweeps vary the external rate $p$ or connectivity
$C$ to change the operating state, as specified in Methods.

\begin{table}[h]\centering\small
\caption{Jansen--Rit parameters \cite{jansen1995}.}
\label{tab:params}
\begin{tabular}{@{}lll@{}}
\toprule
Symbol & Meaning & Value\\
\midrule
$A,\,B$ & excit./inhib.\ PSP gain & $3.25,\ 22$ mV\\
$a,\,b$ & synaptic rates & $100,\ 50$ s$^{-1}$\\
$v_0,\,e_0,\,\rho$ & sigmoid inflection, max rate$/2$, slope & $6$ mV, $2.5$ s$^{-1}$, $0.56$ mV$^{-1}$\\
$C$ & global connectivity & $135$\\
$C_1,C_2,C_3,C_4$ & couplings & $C,\,0.8C,\,0.25C,\,0.25C$\\
$p$ & external input (bifurcation knob) & $0$--$400$ Hz\\
\bottomrule
\end{tabular}
\end{table}

The alpha limit cycle lies between Hopf points near $p=90$ and
$316$~s$^{-1}$. It coexists with a lower stable fixed point up to the
saddle-node on an invariant circle (SNIC, $p\approx114$), then with a
large-amplitude cycle up to its saddle-node near $p=137$. The SNIC
terminates the large-amplitude branch, not the alpha cycle.


\subsection{Static rate with dynamic synapses}
A second-order synapse with time constant $\tau_s$ obeys
\begin{equation}
r(t)\simeq \frac{F(Q(t))}{\tau_m},\qquad
\tau_s^2\ddot r_{\rm syn}+2\tau_s\dot r_{\rm syn}+r_{\rm syn}=r(t),
\label{eq:slow_synapse}
\end{equation}
where $r_{\rm syn}$ is the filtered rate entering the recurrent drive in $Q$, with the original
coupling normalization retained. This is the slow-synapse limit: shortening the
effective membrane time makes the transfer appear static while the network retains
its synaptic dynamics.

\subsection{Static and dynamic MPR mixing coefficients}\label{app:mpr-coefficients}
In this subsection and the derivation that follows, time is measured in units
of $\tau_m$, rate in units of $1/\tau_m$, and $\omega$ denotes angular frequency
in membrane-time units. Thus $r^*=F(Q)$ and physical angular frequency is
$\omega/\tau_m$. Consider an uncoupled MPR population at drive $Q$, with a weak
current $I_{\rm ac}\cos\omega t$.
To second order in $I_{\rm ac}$, the carrier-averaged shift of $r$ is $I_{\rm ac}^2\kappa_{\rm DC}(Q,\omega)$, where
\begin{equation}
\kappa_{\rm DC}(Q,\omega)=\frac{r^*\,[\omega^2+3d^2-\nu^2]}{(d^2+\nu^2)\{[d^2+\nu^2-\omega^2]^2+4d^2\omega^2\}},
\qquad r^*=F(Q),\quad d=\frac{\Delta}{\pi r^*},\quad \nu=2\pi r^*.
\label{eq:crossover}
\end{equation}
The physical rate shift is $I_{\rm ac}^2\kappa_{\rm DC}/\tau_m$.
The fixed-point eigenvalues in normalized time are $-d\pm i\nu$: $d$ is the relaxation rate and $\nu$
the ringing angular frequency, both in membrane-time units; their physical
values are $d/\tau_m$ and $\nu/\tau_m$. Expansion about this fixed point gives the limits
\begin{equation}
\kappa_{\rm DC}(Q,\omega)\xrightarrow[\ \omega\to0\ ]{}\tfrac14F''(Q),\qquad
\kappa_{\rm DC}(Q,\omega)\xrightarrow[\ \omega\gg d,\nu\ ]{}\frac{F'(Q)}{2\omega^2}.
\label{eq:crossover_limits}
\end{equation}
In the first limit the population tracks its static curve, whose curvature rectifies
the input. In the second, the voltage excursion is approximately $I_{\rm ac}/\omega$; its
square supplies a slow drive that the population reads through $F'$. Above the
inflection of $F$, these limits have opposite signs and the coefficient vanishes at
$\omega_z=\sqrt{\nu^2-3d^2}$. The $300$~Hz and kHz carriers
are on the fast side for the parameters of the pyramidal--interneuron network gamma
(PING) circuit studied here. That classification is
parameter-dependent: a shorter effective membrane time can bring the same physical
carrier into the adiabatic regime.


\subsection{Single-cell rate and dynamic MPR mixing}\label{app:mpr-mixing}
The MPR steady-rate curve $F$ of \S\ref{sec:dynamic}, \ifsupp main-text\fi\ Eq.~\eqref{eq:Fmpr}, averages the single-neuron rate $f(q)=\sqrt q/\pi$ for
$q>0$, and $f(q)=0$ otherwise, over the excitability distribution. Here $\eta$ is
the neuron's excitability and $q=\eta+I$ its constant total drive.

The coefficient in Eq.~\eqref{eq:crossover} follows by expanding about the
fixed point $(r^*,v^*)$, with $v^*=-d/2$, in membrane-time units.
The Jacobian and unit-current phasor are
\[
\mathsf J=\begin{pmatrix}2v^*&2r^*\\-2\pi^2r^*&2v^*\end{pmatrix},\qquad
\mathbf h=(i\omega\mathsf I_2-\mathsf J)^{-1}\binom{0}{1},
\]
where $\mathsf I_2$ is the identity matrix and the first-order response is
$I_{\rm ac}\operatorname{Re}(\mathbf h e^{i\omega t})$.
Writing $h_r,h_v$ for its rate and voltage components, with an overbar denoting
complex conjugation and $\operatorname{Re}$ the real part, the carrier-averaged
quadratic source per $I_{\rm ac}^2$ is
\[
\mathbf b_2=\binom{\operatorname{Re}(h_r\overline{h_v})}{\big(|h_v|^2-\pi^2|h_r|^2\big)/2}.
\]
The mean second-order correction is $-I_{\rm ac}^2\mathsf J^{-1}\mathbf b_2$;
its rate component gives Eq.~\eqref{eq:crossover}.
